\section{CEREBRUM: Academic Source Repository (Batch Index)}

This document provides a conso\-lidated list of found\-ational papers, books, and datasets cited in the CEREBRUM framework development, plus the full index of CEREBRUM Technical Reports (PAPER\_001–037). Each external entry is hyperlinked to its official digital record or Open Access PDF.

\begin{itemize}
\item \textbf{hebb1949}: Hebb, D. (1949). The organ\-ization of behavior. \href{https://archive.org/details/organizationofbe00hebb}{View Source}
\item \textbf{vaswani\-2017\-attention}: Vaswani et al. (2017). Attention is all you need. \href{https://arxiv.org/pdf/1706.03762.pdf}{Download PDF}
\item \textbf{bipoo1998}: Bi \& Poo (1998). Synaptic modif\-ications in neurons. \href{https://www.jneurosci.org/content/jneuro/18/24/10464.full.pdf}{Download PDF}
\item \textbf{traag\-2019\-louvain}: Traag et al. (2019). From Louvain to Leiden. \href{https://arxiv.org/pdf/1810.08473.pdf}{Download PDF}
\item \textbf{blondel\-2008\-louvain}: Blondel et al. (2008). Fast unfolding of communities. \href{https://arxiv.org/pdf/0803.0476.pdf}{Download PDF}
\item \textbf{raghavan\-2007\-lpa}: Raghavan et al. (2007). Near linear time community detection. \href{https://arxiv.org/pdf/0709.2931.pdf}{Download PDF}
\item \textbf{page\-1999\-pagerank}: Page et al. (1999). The PageRank citation ranking. \href{http://ilpubs.stanford.edu:8090/422/1/1999-66.pdf}{Download PDF}
\item \textbf{thompson\-1933\-bayesian}: Thompson, W. (1933). On the likelihood of unknown probability. \href{https://www.gwern.net/docs/statistics/bayes/1933-thompson.pdf}{Download PDF}
\item \textbf{russo\-2018\-thompson}: Russo et al. (2018). A tutorial on Thompson sampling. \href{https://web.stanford.edu/~bvr/pubs/TS_Tutorial.pdf}{Download PDF}
\item \textbf{schon\-emann\-1966\--procrustes}: Schonemann, P. (1966). Normalized solution of Procrustes. \href{https://psychometrika.org/index.php/psychometrika/article/download/10.1007/BF02289451/1056}{Download PDF}
\item \textbf{gower\-2004\-procr\-ustes}: Gower \& Dijkst\-erhuis (2004). Procrustes problems. \href{https://openlibrary.org/books/OL8850117M/Procrustes_problems}{View Source}
\item \textbf{ying\-2019\-gnnex\-plainer}: Ying et al. (2019). Gnnexp\-lainer. \href{https://arxiv.org/pdf/1903.03894.pdf}{Download PDF}
\item \textbf{ribeiro\-2016\-lime}: Ribeiro et al. (2016). "Why should I trust you?" LIME. \href{https://arxiv.org/pdf/1602.04938.pdf}{Download PDF}
\item \textbf{bordes\-2013\-transe}: Bordes et al. (2013). Translating embeddings (TransE). \href{https://papers.nips.cc/paper/5071-translating-embeddings-for-modeling-multi-relational-data.pdf}{Download PDF}
\item \textbf{sun\-2019\-rotate}: Sun et al. (2019). RotatE: Relational Rotation. \href{https://arxiv.org/pdf/1902.10197.pdf}{Download PDF}
\item \textbf{lacroix\-2020\-tntco\-mplex}: Lacroix et al. (2020). Temporal Knowledge Base Completion. \href{https://openreview.net/pdf?id=Hye_S0NKvB}{Download PDF}
\item \textbf{finn\-2017\-maml}: Finn et al. (2017). MAML meta-learning. \href{https://arxiv.org/pdf/1703.03400.pdf}{Download PDF}
\item \textbf{horrocks\-2004\-owl}: Horrocks et al. (2004). The OWL Web Ontology Language. \href{https://www.w3.org/TR/owl-features/}{View Source}
\item \textbf{metaqa2017}: Zhang et al. (2018). Variational reasoning (MetaQA). \href{https://arxiv.org/pdf/1709.04071.pdf}{Download PDF}
\item \textbf{webqsp2016}: Yih et al. (2016). Semantic parse labeling (WebQSP). \href{https://aclanthology.org/P16-2033.pdf}{Download PDF}
\item \textbf{hetionet\-2017\-}: Himmelstein et al. (2017). Systematic integration (Hetionet). \href{https://elifesciences.org/articles/26726.pdf}{Download PDF}
\item \textbf{miller2019-expla\-inability}: Miller, T. (2019). Explanation in AI: Social Sciences. \href{https://arxiv.org/pdf/1706.07269.pdf}{Download PDF}
\item \textbf{samek2017-explainable}: Samek et al. (2017). Explainable AI foundations. \href{https://arxiv.org/pdf/1708.08296.pdf}{Download PDF}
\item \textbf{lundberg\-2017\-shap}: Lundberg \& Lee (2017). A unified approach (SHAP). \href{https://papers.nips.cc/paper/7062-a-unified-approach-to-interpreting-model-predictions.pdf}{Download PDF}
\item \textbf{rosvall\-2008\-infomap}: Rosvall et al. (2008). Maps of random walks (Infomap). \href{https://www.pnas.org/content/pnas/105/4/1118.full.pdf}{Download PDF}
\item \textbf{fortunato\-2010\--community}: Fortunato, S. (2010). Community detection in graphs. \href{https://arxiv.org/pdf/0906.0612.pdf}{Download PDF}
\item \textbf{reimers\-2019\-sente\-ncebert}: Reimers et al. (2019). Sentence-BERT. \href{https://arxiv.org/pdf/1908.10084.pdf}{Download PDF}
\item \textbf{velic\-kovic\-2018\-gat}: Velickovic et al. (2018). Graph Attention Networks (GAT). \href{https://arxiv.org/pdf/1710.10903.pdf}{Download PDF}
\item \textbf{trouillon\-2016\-complex}: Trouillon et al. (2016). Complex embeddings (ComplEx). \href{https://arxiv.org/pdf/1606.06357.pdf}{Download PDF}
\item \textbf{nichol\-2018\-reptile}: Nichol et al. (2018). On first-order meta-learning (Reptile). \href{https://arxiv.org/pdf/1803.02999.pdf}{Download PDF}
\item \textbf{jin\-2020\-renet}: Jin et al. (2020). Recurrent event network (RENet). \href{https://arxiv.org/pdf/2004.13524.pdf}{Download PDF}
\item \textbf{qu2020tdgnn}: Qu et al. (2020). Temporal dynamics-aware GNN. \href{https://arxiv.org/pdf/2002.04873.pdf}{Download PDF}
\item \textbf{dong\-2014\-knowl\-edgevault}: Dong et al. (2014). Knowledge Vault. \href{https://research.google/pubs/knowledge-vault-a-web-scale-infrastructure-for-data-fusion.pdf}{Download PDF}
\item \textbf{bertossi\-2011\-database}: Bertossi, L. (2011). Database incon\-sistency constraints. \href{https://www.morganclaypool.com/doi/abs/10.2200/S00305ED1V01Y201010DTM011}{View Source}
\item \textbf{bloom1970}: Bloom, B. (1970). Bloom Filter. \href{http://crystal.uta.edu/~mcguigan/cse6339/papers/Bloom70.pdf}{Download PDF}
\item \textbf{dieke\-lmann\-2010\-memory}: Diekelmann \& Born (2010). The memory function of sleep. \href{https://www.nature.com/articles/nrn2762.pdf}{Download PDF}
\item \textbf{doan\-2012\-princ\-iples}: Doan et al. (2012). Principles of data integration. \href{https://www.sciencedirect.com/book/9780124158146/principles-of-data-integration}{View Source}
\item \textbf{gentry2009}: Gentry, C. (2009). Fully homomorphic encryption. \href{https://crypto.stanford.edu/craig/craig-thesis.pdf}{Download PDF}
\item \textbf{wu2012probase}: Wu et al. (2012). Probase: Probab\-ilistic taxonomy. \href{https://www.cs.uwyo.edu/~ruben/cs5011/probase.pdf}{Download PDF}
\item \textbf{carlson\-2010\-nell}: Carlson et al. (2010). Never-ending language learning (NELL). \href{http://www.cs.cmu.edu/~acarlson/papers/carlson-aaai10.pdf}{Download PDF}
\item \textbf{yang\-2013\-bigclam}: Yang \& Leskovec (2013). Overlapping community (BIGCLAM). \href{https://arxiv.org/pdf/1205.6228.pdf}{Download PDF}
\item \textbf{coscia\-2012\-demon}: Coscia et al. (2012). Local-first overlapping (DEMON). \href{https://arxiv.org/pdf/1206.0629.pdf}{Download PDF}
\item \textbf{lin\-2015\-ttranse}: Lin et al. (2015). Knowledge graph completion (TTransE). \href{https://thunlp.org/~lzy/publications/aaai2015_transr.pdf}{Download PDF}
\item \textbf{sun2017hyte}: Sun et al. (2017). Temporally Aware Embedding (HyTE). \href{https://www.researchgate.net/profile/Shibiao_Sun/publication/320448102_Hyte_Joint_embedding_of_entities_relations_and_time/links/5a9e59d9a6fdcc546ad881f3/Hyte-Joint-embedding-of-entities-relations-and-time.pdf}{Download PDF}
\item \textbf{saxena\-2020\-improve}: Saxena et al. (2020). Improving multi-hop QA (EmbedKGQA). \href{https://arxiv.org/pdf/1909.11855.pdf}{Download PDF}
\item \textbf{hamilton\-2017\-graphsage}: Hamilton, W., Ying, Z., \& Leskovec, J. (2017). Inductive Repres\-entation Learning on Large Graphs. NeurIPS. \href{https://arxiv.org/pdf/1706.02216.pdf}{Download PDF}
\item \textbf{heinlein\-1949\-gulf}: Heinlein, R.A. (1949). Gulf. \textit{Astounding Science Fiction}, Vol. 44, No. 4. The fictional origin of SpeedTalk — a constructed language assigning single phonemes to primitive concepts, inspiring the phonemic compression design in Paper 021. \href{https://archive.org/details/Astounding_v44n04_1949-12}{View Source}
\item \textbf{zhu\-2025\-loooplm}: Zhu, R.-J., Wang, Z., Hua, K., et al. (2025). Scaling Latent Reasoning via Looped Language Models (Ouro). arXiv:2510.25741. ByteDance Seed / UC Santa Cruz et al. Demons\-trates that applying the same transformer stack T times yields drama\-tically better reasoning on hard inputs without increasing parameter count; introduces the adaptive exit gate $\lambda$\_t (ideal conti\-nuation probability) that prevents both under\-thinking and overt\-hinking. Directly inspired CEREBRUM Phase 70 LoopedBeamTraversal. \href{https://arxiv.org/pdf/2510.25741}{Download PDF}
\item \textbf{bengio\-2025\-soliton}: Bengio, Y. et al. (2025). Consci\-ousness as a Soliton, Not a Process: Identity, Memory, and the Hard Problem in Coherence Field Theory. r/CoherencePhysics Preprint, UCFT 2025. Proposes that identity and consc\-iousness are best modeled as stable, self-reinforcing waves (solitons) rather than sequential processes. Framing adopted in CEREBRUM Phase 69 for the \texttt{soliton\_index} — a localized prior that consi\-stently yields low prediction error is soliton-like, maintaining its shape through propagation. \href{https://www.reddit.com/r/CoherencePhysics/comments/1sh749z/consciousness_as_a_soliton_not_a_process_identity/}{View Source}
\item \textbf{friston\-2005\-theory}: Friston, K. (2005). A theory of cortical responses. \textit{Philos\-ophical Transa\-ctions of the Royal Society B}, 360(1456), 815–836. Predictive coding / free energy principle — found\-ational for Phase 69 PredictiveCodingEngine. \href{https://doi.org/10.1098/rstb.2005.1622}{View Source}
\item \textbf{rao\-1999\-predi\-ctive}: Rao, R.P.N. \& Ballard, D.H. (1999). Predictive coding in the visual cortex: a functional inter\-pretation of some extra-classical receptive-field effects. \textit{Nature Neuros\-cience}, 2(1), 79–87. \href{https://www.nature.com/articles/nn0199_79.pdf}{Download PDF}
\item \textbf{bottou\-2010\-large}: Bottou, L. (2010). Large-scale machine learning with stochastic gradient descent. \textit{COMPSTAT}, 177–186. Online SGD basis for AutoApprover (Phase 71). \href{https://link.springer.com/content/pdf/10.1007/978-3-7908-2604-3_16.pdf}{Download PDF}
\item \textbf{denzin\-1978\-research}: Denzin, N.K. (1978). \textit{The research act: A theoretical intro\-duction to socio\-logical methods.} McGraw-Hill. Triang\-ulation methodology — found\-ational framing for Phase 72 Triang\-ulationEngine.
\item \textbf{pearl\-2000\-causality}: Pearl, J. (2000). \textit{Causality: Models, reasoning, and inference.} Cambridge University Press. Noisy-OR model — basis for Contra\-dictionResolver (Phase 73). \href{https://doi.org/10.1017/CBO9780511803161}{View Source}
\item \textbf{bastian\-2009\-gephi}: Bastian, M., Heymann, S., Jacomy, M. (2009). Gephi: An Open Source Software for Exploring and Manipu\-lating Networks. \textit{ICWSM}. Spatial graph layout inspiration for the Fibonacci sphere design in Phase 92. \href{https://gephi.org/}{View Source}
\item \textbf{fibonacci\_sphere}: Álvarez, R. (2010). How to evenly distribute points on a sphere more effectively than the canonical Fibonacci Lattice. Algorithm used for community center placement in \texttt{setup\_graph\_layout.py} (Phase 92). \href{https://arxiv.org/abs/0912.4540}{View Source}

\end{itemize}
---

\subsection{CEREBRUM Technical Reports (PAPER\_001–037)}

\begin{center}
{\footnotesize
\begin{tabularx}{\columnwidth}{|>{\raggedright\arraybackslash}X|>{\raggedright\arraybackslash}X|>{\raggedright\arraybackslash}X|>{\raggedright\arraybackslash}X|}
\hline
Paper & Phase & Topic & File \\ \hline
PAPER\_001 & 5 & DSCF/TSC Community Fusion & \texttt{PAPER\_001\_DSCF\_TSC.md} \\ \hline
PAPER\_002 & 5 & Community-Structured Attention (CSA) & \texttt{PAPER\_002\_CSA.md} \\ \hline
PAPER\_003 & 12 & Bridge Twin Nodes & \texttt{PAPER\_003\_BRIDGE\_TWINS.md} \\ \hline
PAPER\_004 & 13 & STDP Causal Edge Inference & \texttt{PAPER\_004\_STDP\_CAUSAL.md} \\ \hline
PAPER\_005 & 9 & Holographic Federated Index & \texttt{PAPER\_005\_HOLOGRAPHIC\_INDEXING.md} \\ \hline
PAPER\_006 & 19 & Bayesian Beam Search & \texttt{PAPER\_006\_BAYESIAN\_BEAM.md} \\ \hline
PAPER\_007 & 15 & REM Cycle Maintenance & \texttt{PAPER\_007\_REM\_CYCLE.md} \\ \hline
PAPER\_008 & 18 & Cross-Modal Signal Alignment & \texttt{PAPER\_008\_SIGNAL\_ENCODER.md} \\ \hline
PAPER\_009 & 18 & THALAMUS Ingestion Pipeline & \texttt{PAPER\_009\_THALAMUS.md} \\ \hline
PAPER\_010 & 20 & Inference Validation & \texttt{PAPER\_010\_INFERENCE\_VALIDATION.md} \\ \hline
PAPER\_011 & 20 & Contra\-diction Materi\-alization & \texttt{PAPER\_011\_CONTRADICTION.md} \\ \hline
PAPER\_012 & 20 & Glass-Box Reasoning Studio & \texttt{PAPER\_012\_REASONING\_STUDIO.md} \\ \hline
PAPER\_013 & 18 & Streaming Engine & \texttt{PAPER\_013\_STREAMING\_ENGINE.md} \\ \hline
PAPER\_014 & 20 & Metaco\-gnitive Verifi\-cation & \texttt{PAPER\_014\_INSIGHT\_ENGINE.md} \\ \hline
PAPER\_015 & 22 & Algorithmic Depth & \texttt{PAPER\_015\_ALGORITHMIC\_DEPTH.md} \\ \hline
PAPER\_016 & 21 & Production Hardening (Holes 1-12) & \texttt{PAPER\_016\_PRODUCTION\_HARDENING.md} \\ \hline
PAPER\_017 & 65 & Hypothesis Materi\-alization & \texttt{PAPER\_017\_HYPOTHESIS\_MATERIALIZATION.md} \\ \hline
PAPER\_018 & 55 & Engram-Steered Traversal & \texttt{PAPER\_018\_ENGRAM\_TRAVERSAL.md} \\ \hline
PAPER\_019 & 57 & Temporal Calibration & \texttt{PAPER\_019\_TEMPORAL\_CALIBRATOR.md} \\ \hline
PAPER\_020 & 60 & MACH Consensus Hierarchies & \texttt{PAPER\_020\_MACH.md} \\ \hline
PAPER\_021 & 58 & SpeedTalk Phonemic Compression & \texttt{PAPER\_021\_SPEEDTALK\_COMPRESSION.md} \\ \hline
PAPER\_022 & 70 & Looped Beam Traversal & \texttt{PAPER\_022\_LOOPED\_TRAVERSAL.md} \\ \hline
PAPER\_023 & 69 & PredictiveCodingEngine & \texttt{PAPER\_023\_PREDICTIVE\_CODING.md} \\ \hline
PAPER\_024 & 71 & AutoApprover & \texttt{PAPER\_024\_AUTOAPPROVER.md} \\ \hline
PAPER\_025 & 72 & Triang\-ulationEngine & \texttt{PAPER\_025\_TRIANGULATION.md} \\ \hline
PAPER\_026 & 73 & DiscoveryCalibrator + Contra\-dictionResolver & \texttt{PAPER\_026\_DISCOVERY\_CALIBRATION.md} \\ \hline
PAPER\_027 & 74 & AutonomousDiscoveryLoop & \texttt{PAPER\_027\_AUTONOMOUS\_LOOP.md} \\ \hline
PAPER\_028 & 75+78 & Studio v2 Dashboard + Provenance Panel & \texttt{PAPER\_028\_STUDIO\_V2.md} \\ \hline
PAPER\_029 & 76 & ProvenanceLedger & \texttt{PAPER\_029\_PROVENANCE\_LEDGER.md} \\ \hline
PAPER\_030 & 77 & Feature Impact Benchmark & \texttt{PAPER\_030\_FEATURE\_IMPACT.md} \\ \hline
PAPER\_031 & 79 & Loop-Provenance Recovery & \texttt{PAPER\_031\_LOOP\_PROVENANCE\_RECOVERY.md} \\ \hline
PAPER\_032 & 80 & GraphAdapter remove\_edge Protocol & \texttt{PAPER\_032\_GRAPH\_ADAPTER\_PROTOCOL.md} \\ \hline
PAPER\_033 & 81 & GraphSnapshot Persistence & \texttt{PAPER\_033\_GRAPH\_SNAPSHOT.md} \\ \hline
PAPER\_034 & 82 & Adaptive Loop Tuning & \texttt{PAPER\_034\_ADAPTIVE\_LOOP\_TUNING.md} \\ \hline
PAPER\_035 & 83 & UE5 Neural Visual\-ization Bridge & \texttt{PAPER\_035\_UE5\_NEURAL\_VISUALIZATION.md} \\ \hline
PAPER\_036 & 150 & The Cingulate Engine: Reasoning Verifi\-cation & \texttt{PAPER\_036\_CINGULATE\_ENGINE.md} \\ \hline
PAPER\_037 & 167 & GraphProfiler and STRB: Autonomous Strategy & \texttt{PAPER\_037\_GRAPHPROFILER\_STRB.md} \\ \hline
PAPER\_100 & — & Conclusion: The CEREBRUM Paradigm & \texttt{PAPER\_100\_CONCLUSION.md} \\ \hline
\end{tabularx}
}
\end{center}

---
\textbf{Reviewed on}: May 5, 2026 for version v2.52.0
